\section*{Hoofdstuk \ref{ch:b2:limb-organogenesis} --- Organogenese: de bouw van de ledemaat van tetrapoden}

\begin{solution}{exo:b2:limb-organogenesis:1}
Proximodistaal: de \omterm{def:b2:limb-organogenesis:bud}{apicale ectodermale richel} langs de top, FGF's.
Anteroposterior: de \omterm{def:b2:limb-organogenesis:bud}{zone van polariserende activiteit} aan de achterrand,
\omterm{prop:b2:limb-organogenesis:zpa}{Sonic hedgehog}. Dorsoventraal: het dorsale ectoderm, Wnt7a.
\end{solution}

\begin{solution}{exo:b2:limb-organogenesis:2}
Stylopodium: opperarmbeen, dijbeen (Hoxa/d 9--10). Zeugopodium: spaakbeen
en ellepijp, scheenbeen en kuitbeen (Hox 11). Autopodium: pols en hand,
enkel en voet (Hox 13).
\end{solution}

\begin{solution}{exo:b2:limb-organogenesis:3}
Ledemaatveld gespecificeerd (Tbx5); de knop groeit uit onder het FGF van
de richel; de ZPA en het dorsale ectoderm geven hem een patroon; cellen die
de \omterm{prop:b2:limb-organogenesis:aer}{progressiezone} verlaten, condenseren tot stralen; kraakbeen
differentieert; gewrichten ontstaan waar kraakbeen wordt onderdrukt;
bloedvaten dringen binnen en bot vervangt kraakbeen; interdigitale cellen
sterven; spiervoorlopers uit de somieten migreren binnen, versmelten en
hechten zich via pezen.
\end{solution}

\begin{solution}{exo:b2:limb-organogenesis:4}
(a) Een stompje met weinig of geen opperarmbeen; (b) opperarmbeen, spaakbeen
en ellepijp, maar geen hand; (c) een volledige ledemaat --- FGF vervangt de
richel.
\end{solution}

\begin{solution}{exo:b2:limb-organogenesis:5}
$x = 40\ln 2 = \qty{28}{\micro m}$, $40\ln 5 = \qty{64}{\micro m}$,
$40\ln 20 = \qty{120}{\micro m}$: territoria van 28, 36 en
\qty{56}{\micro m}.
\end{solution}

\begin{solution}{exo:b2:limb-organogenesis:6}
(a) 4-3-2-2-3-4, een volledige spiegelreeks. (b) Een zwakkere bron geeft een
lagere piek, zodat alleen de vinger met de lage drempel wordt toegevoegd:
2-2-3-4. (c) Een korte blootstelling specificeert eveneens alleen de meest
anterieure extra vinger: 2-2-3-4 --- cellen integreren de concentratie over
de tijd.
\end{solution}

\begin{solution}{exo:b2:limb-organogenesis:7}
$5000\times 2^{8} = 1.3\times 10^{6}$. Om $4\times 10^{6}$ te bereiken zijn
$\log_2 800 = 9.6$ verdubbelingen nodig, zodat de cyclus gemiddeld ongeveer
\qty{10}{h} moet duren, of de telling de spiervoorlopers moet omvatten die
vanuit de somieten binnenmigreren; celdood zou de discrepantie groter maken,
niet kleiner.
\end{solution}

\begin{solution}{exo:b2:limb-organogenesis:8}
Interdigitale cellen van een kuiken ontvangen een BMP-signaal en sterven;
interdigitale cellen van een eend brengen een BMP-remmer tot expressie en
overleven, zodat het vlies blijft bestaan. De remmer op een kraal redt het
kuikenvlies: een voet met zwemvliezen. Celdood volgens schema in een
transplantaat betekent dat de cellen al vóór het verwijderen waren
geïnstrueerd --- de beslissing wordt vroeg genomen en later, autonoom,
uitgevoerd.
\end{solution}

\begin{solution}{exo:b2:limb-organogenesis:9}
Zonder Hoxa13 en Hoxd13: geen autopodium --- de voorpoot eindigt bij de pols.
Zonder Hoxa11 en Hoxd11: een sterk verkort spaakbeen en ellepijp, met een
hand die vrijwel rechtstreeks aan het opperarmbeen vastzit. Elk segment
heeft zijn eigen Hox-paar nodig; de domeinen zijn niet alleen merkers, maar
instructies.
\end{solution}

\begin{solution}{exo:b2:limb-organogenesis:10}
FGF uit de richel houdt Shh in de ZPA aan; Shh houdt, via een doorgeefluik
in het mesoderm, FGF in de richel aan. Het signaal van elk centrum bevestigt
dus dat het andere werkt, zodat een ledemaat nooit groeit zonder een patroon
te krijgen of een patroon krijgt zonder te groeien, en een tijdelijk verlies
van een van beide door het andere wordt hersteld. Als Shh FGF niet in stand
kon houden, zou de richel binnen een dag in regressie gaan en zou de uitgroei
stoppen bij het segment dat dan was bereikt: een afgeknotte ledemaat. De lus
van de bevalling wordt beëindigd door de gebeurtenis die ze aandrijft ---
de geboorte neemt de prikkel weg; de lus van de ledemaat wordt beëindigd
door differentiatie --- het mesoderm geeft niets meer door wanneer de
ledemaat af is.
\end{solution}

\begin{solution}{exo:b2:limb-organogenesis:11}
$e^{4r} = 8$: $r = \ln 8/4 = \qty{0.52}{d^{-1}}$. Groei: een langere,
snellere proliferatie van het vingerkraakbeen (een ledemaatenhancer die de
expressie van een groeigen in de vingers verhoogt). Dood: overleving van
het interdigitale weefsel door expressie van een BMP-remmer, die het
membraan behoudt.
\end{solution}

\begin{solution}{exo:b2:limb-organogenesis:12}
Signalen worden hergebruikt omdat een werkende module van receptor en
signaaltransductie goedkoop opnieuw in te zetten is: wat tussen weefsels
verandert, is niet het signaal maar de verzameling genen die de ontvangende
cellen competent zijn aan te schakelen, zodat hetzelfde
concentratieprofiel op de ene plaats motorneuron betekent en op een andere
vinger 4. Een \omterm{def:b2:genome-diversification:mutation}{mutatie} in zo'n gen treft daardoor veel organen tegelijk ---
\omterm{def:b2:genome-diversification:mutation}{mutaties} van Shh geven holoprosencefalie en ledemaatafwijkingen samen ---
tenzij ze in een weefselspecifieke enhancer ligt, in welk geval ze maar één
orgaan verandert, en zo verloopt de meeste evolutionaire verandering in
deze genen.
\end{solution}

\begin{solution}{pb:b2:limb-organogenesis:1}
\textbf{1.} $96/12 = 8$ verdubbelingen: $2000\times 256 = 5.1\times 10^{5}$
cellen.
\textbf{2.} $1000/256 \approx 4$: celvergroting en kraakbeenmatrix.
\textbf{3.} Lengte $\times 10$; een cilinder met constante straal zou
$\times 10$ in volume geven; de extra factor honderd betekent dat ook de
straal ongeveer tienvoudig groeide --- de knop verbreedde zich tot een
peddel terwijl hij langer werd.
\textbf{4.} Na 3 dagen $300/400 = \qty{75}{\%}$; na 7 dagen $300/4000 =
\qty{7.5}{\%}$.
\textbf{5.} Alleen de top staat onder het signaal; naarmate de top oprukt,
verlaten de achtergebleven cellen de zone in de volgorde waarin ze zijn
gemaakt, zodat proximale structuren eerst worden gespecificeerd en distale
als laatste.
\textbf{6.} Het zeugopodium (spaakbeen en ellepijp), gespecificeerd tussen
3.5 en 4 dagen.
\textbf{7.} Stylopodium na 3.5 dagen, zeugopodium na 4.0 dagen, autopodium
na 4.5 dagen (vingertoppen later).
\textbf{8.} Een halve dag: één celcyclus.
\textbf{9.} Elke cyclus in de \omterm{prop:b2:limb-organogenesis:aer}{progressiezone} schuift de Hox-code op (9--10,
dan 11, dan 13); cellen die na één, twee of drie cycli vertrekken, dragen
de code van het stylopodium, het zeugopodium of het autopodium.
\textbf{10.} Een volledige vleugel: FGF is wat de richel leverde.
\textbf{11.} Een tweede uitgroei vanuit het dorsale oppervlak: verdubbelde
distale structuren, een overtallige ledemaat.
\textbf{12.} Het mesoderm geeft Shh niet meer door aan de richel, de
expressie van Shh stopt, de richel gaat in regressie, en cellen
differentiëren in plaats van te delen: de lus die de groei onderhield, is
verbroken.
\textbf{13.} $k = \ln 2/1740 = \qty{4.0e-4}{s^{-1}}$; $\lambda =
\sqrt{10^{-12}/4\times 10^{-4}} = \qty{50}{\micro m}$.
\textbf{14.} $x = 50\ln(1/0.6) = 26$, $50\ln 4 = 69$, $50\ln 12.5 =
\qty{126}{\micro m}$: territoria van 26, 43 en \qty{57}{\micro m}.
\textbf{15.} $600 - 126 = \qty{474}{\micro m}$.
\textbf{16.} Elke grens schuift met $50\ln 2 = \qty{35}{\micro
m}$ naar buiten: 61, 104, 161; de buitengrens van het territorium van
vinger 2 komt op \qty{161}{\micro m} in plaats van 126, genoeg voor een
extra anterieure vinger 2.
\textbf{17.} Elke reeks heeft \qty{126}{\micro m} nodig: minstens
\qty{252}{\micro m}; de knop van \qty{600}{\micro m} heeft ruimte, met
\qty{348}{\micro m} vingervrij weefsel in het midden.
\textbf{18.} Bij \qty{200}{\micro m} overlappen de twee gradiënten en tellen
ze op: de middelste cellen zien genoeg voor vinger 3 en de territoria van
vinger 2 verdwijnen --- 4-3-4 of 4-3-3-4, minder vingers dan een volledige
spiegelreeks.
\textbf{19.} Een kwart van de cellen geeft $c_0/4$: geen enkele cel bereikt
de drempel van vinger 4 of 3 ($0.25 c_0$ aan de rand zelf), en de drempel
van vinger 2 wordt overschreden tot op $50\ln(0.25/0.08) =
\qty{57}{\micro m}$: één extra vinger 2.
\textbf{20.} $3\times 300\times 300\times 200 = 5.4\times 10^{7}$
\unit{\micro m^3}: $5.4\times 10^{4}$ cellen; over \qty{86400}{s}
$0.6$ cellen per seconde.
\textbf{21.} Kraal met BMP in het eendenvlies: het vlies sterft af en de
tenen komen los; remmer in het kuikenvlies: het vlies overleeft, een voet
met zwemvliezen.
\textbf{22.} $e^{0.5\times 4} = e^{2} = 7.4$. De tweede verandering is het
overleven van het interdigitale membraan dankzij een BMP-remmer.
\textbf{23.} De late fase van Hox 13-expressie die een autopodium specificeert,
en het blijven bestaan van de richel: de vinplooi van de vis maakt in plaats
daarvan vinstralen. De vin van \emph{Tiktaalik} bevat polsachtige botjes ---
een vin die een ledemaat begint te worden.
\textbf{24.} $5.4\times 10^{4}/1000 = 54$: $\log_2 54 = 5.8$ verdubbelingen,
ongeveer \qty{69}{h}, drie dagen om te maken wat één dag verwijdert.
\textbf{25.} Acht verdubbelingen; vingergrenzen bij 26, 69 en
\qty{126}{\micro m}; een verdubbeling in spiegelbeeld vraagt minstens
\qty{252}{\micro m}.
\end{solution}
